\section{Background}
\label{sec:background}

This section fixes notation for the four systems Papilot sits between: the
\TeX{} build, the PDF format, Git, and SyncTeX.

\subsection{The \LaTeX{} build as a fixed point}

\LaTeX{} resolves cross-references, citations and the table of contents by
reading auxiliary data written by the previous run~\cite{lamport1994}. Let
$x_k$ be the auxiliary state after run $k$ and $F$ the map computed by one run,
so that
\begin{equation}
  x_{k+1} = F(x_k), \qquad x_0 = \varnothing .
  \label{eq:fixed-point}
\end{equation}
A build tool such as \texttt{latexmk} repeats runs until a fixed point
$x^{\star} = F(x^{\star})$ is reached or a limit of $K$ runs is hit (five by
default). A single run suffices when $F(\varnothing) = \varnothing$, which is
the case only for documents without references. A typical paper converges in
two or three runs: the first records labels, the second resolves them, and a
third is needed only when resolving them moved a page break.

\begin{remark}
$F$ need not be a contraction. A reference whose resolution lengthens a line
can push a label onto the next page, whose page number shortens the line again;
the sequence $x_k$ then alternates between two states and the run limit ends
the build. Such documents are rare, but they are why a build key must include
the build settings and not only the sources (\cref{sec:build}).
\end{remark}

\subsection{Text in a PDF}

A PDF page places glyphs with a chain of matrices~\cite{iso32000}. The text
state contains the font size $T_{fs}$, the horizontal scaling $T_h$, the rise
$T_{\text{rise}}$, the character and word spacings $T_c$ and $T_w$, and the
text matrix $T_m$. A glyph is drawn in user space through the text rendering
matrix
\begin{equation}
  T_{rm} =
  \begin{pmatrix}
    T_{fs} T_h & 0 & 0 \\
    0 & T_{fs} & 0 \\
    0 & T_{\text{rise}} & 1
  \end{pmatrix}
  \times T_m \times \mathrm{CTM},
  \label{eq:trm}
\end{equation}
where $\mathrm{CTM}$ is the current transformation matrix. After each glyph
with horizontal advance $w_0$ (in thousandths of a text space unit) and a
positioning adjustment $T_j$, the text matrix moves by
\begin{equation}
  t_x = \Bigl( \bigl(w_0 - \tfrac{T_j}{1000}\bigr) T_{fs} + T_c + T_w \Bigr) T_h ,
  \qquad
  T_m \leftarrow
  \begin{pmatrix} 1 & 0 & 0 \\ 0 & 1 & 0 \\ t_x & 0 & 1 \end{pmatrix} T_m ,
  \label{eq:advance}
\end{equation}
with $T_w$ applied only to the single-byte space character. \TeX{} rarely uses
$T_w$; it positions words with explicit adjustments $T_j$, which is why a PDF
produced from \TeX{} often contains no space characters at all.

\subsection{Git worktrees}

A Git repository is a store of immutable objects (blobs, trees and commits)
addressed by their hashes, together with a set of references~\cite{chacon2014}.
A \emph{worktree} is a checkout with its own working files $W_i$, index $I_i$
and head $\mathrm{HEAD}_i$, sharing the object store with every other worktree
of the same repository:
\begin{equation}
  \mathcal{R} = \bigl(\mathcal{O},\ \mathrm{refs}\bigr), \qquad
  w_i = \bigl(W_i,\ I_i,\ \mathrm{HEAD}_i\bigr), \quad i = 0, 1, \dots, t .
\end{equation}
The manuscript checkout is $w_0$; each task $i \ge 1$ gets its own $w_i$ on its
own branch. Because objects are shared, creating a worktree costs only the
working files, and because indexes are separate, a task can stage, diff and
build without touching $w_0$.

\subsection{SyncTeX}

SyncTeX records, while \TeX{} typesets, where each box, glue and kern came from:
an input file, a line and a column~\cite{laurens2008}. Positions are recorded
in \TeX{}'s scaled points. With
\begin{equation}
  1\,\mathrm{pt} = \tfrac{1}{72.27}\,\mathrm{in}, \qquad
  1\,\mathrm{bp} = \tfrac{1}{72}\,\mathrm{in}, \qquad
  1\,\mathrm{sp} = 2^{-16}\,\mathrm{pt},
\end{equation}
a horizontal position $h$ in scaled points becomes PDF points as
\begin{equation}
  x = \bigl(h + h_{\text{origin}}\bigr) \cdot 2^{-16} \cdot \frac{72}{72.27},
  \label{eq:sp-to-bp}
\end{equation}
where $h_{\text{origin}}$ is the one-inch offset of \TeX{}'s origin from the
page edge, expressed in scaled points. A reverse query (``edit'') returns the
record nearest to a point on the page; a forward query (``view'') returns the
boxes produced by a source line. Papilot uses only reverse queries.
